% source: erdos-872/researcher-57-pro-round15-bonferroni4-PROVED-L-le-0.19n.md
% source: erdos-872/r57_bonferroni4_audit_and_repair.md
% source: erdos-872/researcher-60-pro-R57-repair-PROVED-theorems-2-1-and-4-1.md
% source: erdos-872/lean/erdos_872_core/RequestProject/Round15Bonferroni4/Target.lean

\section{The Main Upper Bound}\label{sec:main-upper}

In this section we prove the paper's headline upper bound
\[
  \L(n) \le \left(\frac{\Wfour}{2}+o(1)\right)n < 0.19n.
\]
The strategy is \emph{strategy-specific} but not $\sigma^*$-dependent: Shortener
plays the smallest legal odd prime for a long initial prefix of
\[
  K = \left\lfloor \frac{(1-\eps)n}{2\log n}\right\rfloor
\]
turns.  Let
\[
  q_1 < q_2 < \cdots < q_K
\]
be the odd primes she plays during that prefix.  The argument has four layers:
\begin{enumerate}
  \item a local density law for the counting function of the $q_j$;
  \item a monotone envelope and inversion producing a comparison sequence
    $b_j \ge q_j$ with explicit moment limits;
  \item a fourth-order Bonferroni comparison theorem for arbitrary increasing
    odd-prime comparison sequences;
  \item a prime-rounding bridge that replaces the real sequence $b_j$ by actual
    odd primes $p_j$, using the key half-density slack supplied by the
    prime-rounding bridge of \Cref{thm:prime-bridge}.
\end{enumerate}

\subsection{The Round~15 density}

Write
\[
  S_K(X) := \#\{j\le K : q_j \le X\},
  \qquad
  u := \frac{\log X}{\log n}.
\]
On the interval
\[
  u \in \left(\frac{1}{h+1},\frac1h\right],
\]
the Round~15 local theorem identifies the limiting captured-prime density as
\[
  \rhoU(u) = \frac{1}{(\lfloor 1/u \rfloor+1)u} = \frac{1}{(h+1)u}.
\]

\begin{proposition}[Local density away from breakpoints]\label{prop:local-density}
Fix $H\ge 2$.  For each $1\le h\le H-1$, choose a compact interval
\[
  I_h \subset \left(\frac{1}{h+1},\frac1h\right)
\]
bounded away from the endpoints by a positive amount depending only on $H$.
Then uniformly for $u\in I_h$,
\[
  S_K(n^u)
  \ge
  (1-o_H(1))\frac{n^u}{(h+1)u\log n}.
\]
\end{proposition}

\begin{proof}[Proof sketch]
This is the Round~15 local prime-count-per-range lemma.  One compares the count
of legal odd primes up to $n^u$ with the number of primes in the interval
$[n^{u-\eta},n^u]$ for a fixed small $\eta>0$ depending only on the distance
from the breakpoints.  On $I_h$ the local packet ratio is constant and equals
$1/(h+1)$, so the prime-count lower bound becomes
\[
  \frac{\pi(n^u)-\pi(n^{u-\eta})}{h+1}-O_H(1).
\]
The prime number theorem then yields the claimed asymptotic lower bound,
uniformly on each compact interval $I_h$.
\end{proof}

\subsection{Moment constants and the Bonferroni coefficient}

Define
\[
  J_r
  :=
  \frac1{r!}
  \int_{\substack{u_1+\cdots+u_r\le 1\\u_i\in(0,1]}}
  \prod_{i=1}^r \rhoU(u_i)\,du_1\cdots du_r,
\]
and
\[
  \Wfour := 1-J_1+J_2-J_3+J_4.
\]
Numerical evaluation gives
\[
  \Wfour = 0.3794224,
  \qquad
  \frac{\Wfour}{2} = 0.1897112 < 0.19.
\]

For any increasing sequence $b_1\le \cdots \le b_K$ of positive reals, define
the truncated factorial-moment sums
\[
  T_r^{(b)}(n)
  :=
  \sum_{\substack{1\le j_1<\cdots<j_r\le K\\b_{j_1}\cdots b_{j_r}\le n}}
  \frac{1}{b_{j_1}\cdots b_{j_r}}
  \qquad (1\le r\le 4).
\]

\subsection{Monotone envelope and inversion}

The local density from \Cref{prop:local-density} is piecewise-defined, so one
must build a monotone global comparison sequence.  This is exactly the content
of Steps~3--5 of the Round~57 note.

\begin{proposition}[Envelope and inversion]\label{prop:envelope-inversion}
For each fixed $H$ there exists a monotone cumulative lower envelope
$C_{H,n}(X)$ such that:
\begin{enumerate}
  \item $C_{H,n}(X)\le S_K(X)$ for all $2\le X\le n$;
  \item if
    \[
      b_j := \inf\{X\in[2,n]: C_{H,n}(X)\ge j\},
    \]
    then $q_j \le b_j$ for every $1\le j\le K$;
  \item after first letting $n\to\infty$ and then $H\to\infty$,
    \[
      T_r^{(b)}(n)=J_r+o(1)
      \qquad (1\le r\le 4).
    \]
\end{enumerate}
\end{proposition}

\begin{proof}[Proof sketch]
On each compact block $I_h$ from \Cref{prop:local-density}, define a lower
envelope
\[
  A_{n,h}(u) \asymp \frac{n^u}{(h+1)u\log n}.
\]
These functions are increasing for large $n$.  The cumulative envelope
$C_{H,n}(X)$ is obtained by stitching together the blockwise floors
$\lfloor A_{n,h}(\log_n X)\rfloor$ and filling the breakpoint gaps by constant
``flat blocks.''  By construction, $C_{H,n}\le S_K$, so the inverse sequence
$b_j$ satisfies $q_j\le b_j$.

The flat blocks contribute negligible reciprocal mass: each logarithmic cell has
mass $O(1/\log n)$, and the tuples crossing the boundary layers contribute only
$o(1)$ to the first four factorial moments.  On the genuine blocks, the weak
convergence of the empirical measure follows from the same local density law,
and this gives
\[
  T_r^{(b)}(n)=J_r+o(1)
\]
for $1\le r\le 4$.
\end{proof}

\subsection{The prime-sequence Bonferroni comparison theorem}

The comparison theorem must be stated for \emph{prime sequences}, not for
arbitrary reals; the correct statement is the following.

\begin{theorem}[Prime-sequence Bonferroni--4 comparison]\label{thm:bonf-comp}
Let $q_1<\cdots<q_K$ be the actual odd primes played by Shortener during the
Round~15 prefix, and let $p_1<\cdots<p_K$ be any increasing sequence of odd
primes with
\[
  q_j \le p_j \qquad (1\le j\le K).
\]
Assume that for $1\le r\le 4$,
\[
  T_r^{(p)}(n) = \Lambda_r + o(1).
\]
Then
\[
  \L(n)
  \le
  \frac n2\left(1-\Lambda_1+\Lambda_2-\Lambda_3+\Lambda_4+o(1)\right).
\]
\end{theorem}

\begin{proof}
Let $A$ be the final played set.  If the game ends before Shortener completes
the prefix of length $K$, then $\L(n)=o(n)$ and there is nothing to prove, so
assume $q_1,\dots,q_K$ were all played.

Define the odd-part map
\[
  \varphi(x) := \frac{x}{2^{v_2(x)}}.
\]
If $\varphi(x)=\varphi(y)=m$, then $x=2^a m$ and $y=2^b m$.  If $a\neq b$, one
of $x$ and $y$ divides the other, contradicting that $A$ is an antichain.  So
$\varphi$ is injective on $A$.

Since each $q_j\in A$, no other element of $A$ can be divisible by any $q_j$.
Because $q_j$ is odd, divisibility by $q_j$ is preserved under the odd-part
map.  Hence
\[
  \varphi(A)
  \subseteq
  \{m\le n : m \text{ odd and } q_j\nmid m \text{ for all } j\}
  \cup \{q_1,\dots,q_K\}.
\]
Therefore
\[
  \L(n) \le N(q_1,\dots,q_K) + K = N(q_1,\dots,q_K)+o(n),
\]
where $N(q_1,\dots,q_K)$ counts odd integers up to $n$ avoiding all $q_j$.

Now replace the actual sequence $q_j$ by the larger prime sequence $p_j$.  The
count of odd integers avoiding a finite prime set is monotone in each prime, so
\[
  N(q_1,\dots,q_K) \le N(p_1,\dots,p_K).
\]
Applying fourth-order Bonferroni to the divisibility events $p_j \mid m$ over
odd integers $m\le n$ gives
\[
\begin{aligned}
  N(p_1,\dots,p_K)
  \le\;&
  \frac n2
  \left(1-T_1^{(p)}(n)+T_2^{(p)}(n)-T_3^{(p)}(n)+T_4^{(p)}(n)\right) \\
  &\quad + O(M_2+M_3+M_4),
\end{aligned}
\]
where $M_r$ is the number of odd integers $\le n$ with at least $r$ distinct
prime divisors from the set $\{p_j\}$.  The standard squarefree divisor-count
bound gives
\[
  M_r = O_r\!\left(\frac{n(\log\log n)^{r-1}}{\log n}\right)=o(n)
  \qquad (2\le r\le 4),
\]
so the error is $o(n)$.  Substituting $T_r^{(p)}(n)=\Lambda_r+o(1)$ yields the
stated inequality.
\end{proof}

\subsection{The prime-rounding bridge}

It remains to convert the real comparison sequence $b_j$ from
\Cref{prop:envelope-inversion} into an actual increasing odd-prime sequence.
The naive claim
\[
  p_j/b_j = 1 + O(1/\log b_j)
\]
uniformly in $j$ is \emph{not} justified by the prime number theorem alone and
must not be used.  The correct bridge is the half-density construction of
\Cref{thm:prime-bridge}.

\begin{theorem}[Prime-rounding bridge]\label{thm:prime-bridge}
Let $b_1\le \cdots \le b_K$ be the increasing sequence obtained from the
Round~57 envelope construction.  Then there exists an increasing sequence of odd
primes
\[
  p_1<\cdots<p_K
\]
such that
\[
  q_j \le b_j \le p_j \qquad (1\le j\le K)
\]
and
\[
  \sup_{b_j\le n}\frac{p_j}{b_j}=1+o(1).
\]
Consequently, for each fixed $1\le r\le 4$,
\[
  T_r^{(p)}(n)=T_r^{(b)}(n)+o(1)=J_r+o(1).
\]
\end{theorem}

\begin{proof}
Let
\[
  B_n(X):=\#\{j:b_j\le X\}.
\]
The decisive input is the local-density envelope bound inherited from the
Round~57 construction:
\[
  \frac{dB_n}{dX}
  =
  \frac{\rhoU(u)}{\log n}
  \le
  \frac{1}{2\log X}
\]
uniformly for $X$ in the Round~57 range.  Equivalently, for every fixed
$\Delta>0$,
\[
  B_n((1+\Delta)X)-B_n(X)
  \le
  \left(\frac12+o(1)\right)\frac{\Delta X}{\log X}.
\]
Thus the $b_j$ occupy each short multiplicative interval with density at most
\emph{half} the odd-prime density.

Fix $\delta>0$ and partition $[n^\alpha,n]$ into multiplicative bins
\[
  I_k=[x_k,x_{k+1}),
  \qquad
  x_{k+1}=(1+\delta)x_k.
\]
Assign the $b_j$ that lie in $I_k$ to distinct odd primes in the \emph{next}
bin
\[
  P_k=[x_{k+1},x_{k+2}).
\]
By the prime number theorem,
\[
  \#P_k = (1+o(1))\frac{\delta x_{k+1}}{\log x_k},
\]
while the local-density bound above shows that the number of $b_j$ in $I_k$ is
at most
\[
  \left(\frac12+o(1)\right)\frac{\delta x_k}{\log x_k}.
\]
For large $n$, each next bin therefore contains more than twice as many odd
primes as there are $b_j$ in the preceding bin.  Greedy assignment within each
bin produces an increasing odd-prime sequence with
\[
  b_j \le p_j \le (1+\delta)^2 b_j.
\]
Because the greedy choice is coordinatewise minimal, it is no larger than any
other feasible increasing prime assignment.

Now choose $\delta=\delta(n)\to 0$ along a diagonal sequence.  This gives
\[
  \sup_{b_j\le n}\frac{p_j}{b_j}=1+o(1).
\]
Let $\lambda_n$ denote this supremum.  Then for every $r$-tuple
$I=\{j_1<\cdots<j_r\}$,
\[
  b_I \le p_I \le \lambda_n^r b_I,
  \qquad
  b_I := \prod_{j\in I} b_j,
  \quad
  p_I := \prod_{j\in I} p_j.
\]
Hence
\[
  T_r^{(p)}(n)\le T_r^{(b)}(n),
\]
and also
\[
  T_r^{(p)}(n)\ge \lambda_n^{-r}T_r^{(b)}(n/\lambda_n^r).
\]
Since $\lambda_n\to 1$ and $T_r^{(b)}(n)=J_r+o(1)=O(1)$, the multiplicative
coefficient error is $o(1)$.  The remaining boundary layer
\[
  T_r^{(b)}(n)-T_r^{(b)}(n/\lambda_n^r)
\]
is also $o(1)$, because the same local-density bound gives reciprocal mass
$O(1/\log n)$ for each logarithmic cell, and only $O((\log n)^{r-1})$ cell
patterns contribute near the boundary.  Therefore
\[
  T_r^{(p)}(n)=T_r^{(b)}(n)+o(1)=J_r+o(1)
\]
for $1\le r\le 4$.
\end{proof}

\subsection{The finite Bonferroni--4 theorem}

\begin{theorem}\label{thm:main-upper}
Under the Round~15 strategy $\sigfifteen$,
\[
  \L(n) \le \left(\frac{\Wfour}{2}+o(1)\right)n
  = (0.1897112+o(1))n < 0.19n.
\]
\end{theorem}

\begin{proof}
By \Cref{prop:envelope-inversion}, the inverted sequence $b_j$ satisfies
\[
  T_r^{(b)}(n)=J_r+o(1)
  \qquad (1\le r\le 4).
\]
By \Cref{thm:prime-bridge}, there is an increasing odd-prime sequence $p_j$
with $q_j\le p_j$ and
\[
  T_r^{(p)}(n)=J_r+o(1)
  \qquad (1\le r\le 4).
\]
Applying \Cref{thm:bonf-comp} with $\Lambda_r=J_r$ gives
\[
  \L(n)
  \le
  \frac n2\left(1-J_1+J_2-J_3+J_4+o(1)\right)
  =
  \left(\frac{\Wfour}{2}+o(1)\right)n.
\]
Numerical evaluation yields
\[
  \frac{\Wfour}{2}=0.1897112<0.19,
\]
so the theorem follows.
\end{proof}

\begin{remark}
The Lean file \texttt{Round15Bonferroni4/Target.lean} formalizes the endgame
reduction: once the first four moment terms are componentwise within
$10^{-5}$ of the certified constants $J_1,\dots,J_4$, the strict inequality
$\L(n)<0.19n$ follows eventually.  Theorems~\ref{thm:bonf-comp} and
\ref{thm:prime-bridge} are the rigorous prose ingredients that supply those
moment hypotheses.
\end{remark}
